% !TEX program = xelatex
\documentclass[UTF8,a4paper,11pt]{ctexart}

\usepackage[a4paper,margin=2.05cm,headheight=14pt]{geometry}
\usepackage{amsmath,amssymb,mathtools,bm}
\usepackage{booktabs,longtable,array,tabularx,multirow}
\usepackage{graphicx}
\usepackage{xcolor}
\usepackage{xurl}
\usepackage{hyperref}
\usepackage{seqsplit}
\definecolor{pyblue}{RGB}{36,91,163}
\definecolor{qudared}{RGB}{190,65,55}
\hypersetup{
  colorlinks=true,
  linkcolor=pyblue,
  urlcolor=pyblue,
  pdftitle={PyQCU/QUDA MultiGrid 重新校验与 double MG 修复}
}
\urlstyle{tt}
\Urlmuskip=0mu plus 3mu
\newcommand{\pathref}[1]{\begingroup\Urlmuskip=0mu plus 3mu\path{#1}\endgroup}
\newcommand{\code}[1]{%
  \ifmmode
    \mathtt{\detokenize{#1}}%
  \else
    \begingroup\ttfamily
    \expandafter\seqsplit\expandafter{\detokenize{#1}}%
    \endgroup
  \fi}
\newcolumntype{L}[1]{>{\raggedright\arraybackslash}p{#1}}
\newcolumntype{Y}{>{\raggedright\arraybackslash}X}
\newcommand{\norm}[1]{\left\lVert#1\right\rVert}
\setlength{\parindent}{2em}
\setlength{\parskip}{0.35em}
\emergencystretch=3em

\title{PyQCU 与 QUDA MultiGrid：旧加速比撤回、\texorpdfstring{double MG}{double MG} 修复与重新对照}
\author{PyQCU 工程分析记录}
\date{2026 年 9 月 17 日}

\begin{document}
\maketitle

\begin{abstract}
本文撤回 2026-09-16 报告中未经对齐预算的 c64 三层 $4.3817\times$ 结论，
在相同 gauge、右端项、canonical null vectors、odd--odd MATPC、
停机容差、零初值和五次 median/MAD 口径下重新测量。正式 trace 环境变量
全部关闭；同侧 plain BiCGStab 作为独立参考，不进入 MG 加速比。

重新构建后，QUDA 以 \code{QUDA_PRECISION=12} 同时提供 c64/c128，并启用
\code{QUDA_MULTIGRID_DOUBLE=ON}。为真正运行 double coarse MG，还需修复
double coarse-link 原子累加仍按 \code{int} 固定点解释的问题；否则 level-1
setup 会产生 NaN 或 norm 检查中止。主口径采用
\code{--quda-strategy aligned}，即双侧都是一次递归 V-cycle；
\code{--quda-strategy library} 保留 QUDA 原生 CA-GCR、\code{nu\_post=8} 与
coarse maxiter 16，作为敏感性对照。

当前可信的大格 c64 三层主结果为 PyQCU $0.662549$ s、QUDA $1.318356$ s，
加速比 $1.9898\times$，外迭代数 $14/39$，真残差分别为
$4.90135\times10^{-7}$ 与 $7.32826\times10^{-7}$。c128 三层在
$16^3\cdot16$ 上得到 PyQCU $0.380444$ s、QUDA $1.230048$ s，
加速比 $3.2332\times$，外迭代数 $44/87$。c128 小格三层复测为
$12.3768\times$。c128 三层 $16\cdot32\cdot32\cdot48$ 的全 double
formal cache-hit 点为 PyQCU $2.065148$ s、QUDA $17.245388$ s，
加速比 $8.3507\times$，外迭代数 $20/60$。该点使用 extended
$C=4$ 构建的 asset-identity 缓存，formal 冷 $C=1$ setup 仍未完成。
\end{abstract}

\tableofcontents

\section{旧结果撤回与验收口径}

旧 c64 $16\cdot32\cdot32\cdot48$ 三层结果
$0.653995$ s 对 $2.865627$ s、$4.3817\times$ 不再作为正式结论。
重新检查后确认旧 QUDA 记录使用了与 PyQCU 不同的 smoother 与中间层
coarse solver 预算：PyQCU 每个外层迭代执行一次递归 V-cycle，而旧 QUDA
配置在中间层执行多轮内层迭代。新旧差异不是浮点残差门限或计时边界造成的。

本轮验收标准为：

\begin{enumerate}
  \item 新增性能测试功能默认关闭；正式速度比必须在 trace 环境变量为空时产生；
  \item PyQCU 与 QUDA 使用同一 gauge、RHS、canonical full null vectors、
        mass、block、coarse spin、odd--odd MATPC 和相对真残差门限；
  \item 至少覆盖二、三层，至少覆盖 c64/c128，并扩展到三个格点体积；
  \item 每侧记录 MG steady solve、独立真残差，以及同侧 plain
        BiCGStab 参考时间与真残差；
  \item 对小格点另测 trace-on，明确量化 trace 对性能的影响；
  \item 对 c128 double MG 做真实 setup/solve 验证，不能只以编译成功作证据；
  \item 在硬件能力范围内给出可追溯的正式 JSON，无法运行的条件明确标记。
\end{enumerate}

\section{实现与修复}

\subsection{默认关闭的 trace 与正式闸门}

PyQCU 使用 \code{PYQCU_STRICT_TRACE_FILE}，QUDA 使用
\code{QUDA_MG_TRACE_FILE}。两者未设置时不创建 trace 文件，也不分配
trace 专用 workspace。正式 profile 若检测到任一变量存在，会在启动 worker
前失败；诊断运行必须显式传 \code{--allow-trace}。这样可避免把分层残差
计算、同步和字符串 I/O 混入正式 speedup。

\subsection{BiCGStab 参考与 QUDA 指针修复}

PyQCU 参考调用 \code{applyCloverBistabCgQcu}；QUDA 参考将
\code{inv\_type} 设为 \code{QUDA\_BICGSTAB\_INVERTER}。仅把
\code{inv\_type\_precondition} 设为 INVALID 不够，因为 PyQUDA 的
\code{getClover()} 仍保留 MG 实例指针。QUDA 的 solver factory 会把
BiCGStab 判定为“不支持显式预条件器”并终止。修复后同时设置
\code{inv\_type\_precondition = QUDA\_INVALID\_INVERTER} 与
\code{invert.preconditioner = Pointer("void")}。

参考解与 MG 使用相同 tolerance、maxiter、零初值、误差复算和真残差门限，
但被明确标记为 \code{excluded_from_speedup=true}，不参与
\code{speedup_pyqcu_over_quda}。

\subsection{QUDA double MultiGrid}

上游 QUDA 1.1.0 默认关闭 double recursive MultiGrid。为得到可运行的
c128 coarse 路径，本轮做了以下最小修复：

\begin{enumerate}
  \item CMake 增加 \code{QUDA_MULTIGRID_DOUBLE}，开启时定义
        \code{GPU\_MULTIGRID\_DOUBLE}；
  \item \code{matrix\_tile.cuh} 对 mixed precision gauge accessor
        显式构造目标 \code{complex<T>}；
  \item double coarse-link 的 shared/global atomic accumulation 使用
        \code{double} 存储，而不是把 double 字段按 \code{int} 固定点解释；
  \item c64 仍保留确定性 \code{int} fixed-point 路径，避免回归性能与行为。
\end{enumerate}

关键实现位于：

\begin{itemize}
  \item \pathref{refer/git-rep/quda/include/kernels/coarse_op_kernel.cuh}
  \item \pathref{refer/git-rep/quda/lib/coarse_op.in.cu}
  \item \pathref{refer/git-rep/quda/lib/coarsecoarse_op.hpp}
  \item \pathref{refer/git-rep/quda/lib/staggered_coarse_op.in.cu}
\end{itemize}

统一构建使用：

\begin{verbatim}
QUDA_PRECISION=12
QUDA_RECONSTRUCT=7
QUDA_INTERFACE_QDP=ON
QUDA_QIO=ON / QUDA_QMP=ON
QUDA_MULTIGRID_NVEC_LIST=12,24
QUDA_MULTIGRID_DOUBLE=ON
QUDA_ENABLE_MMA=OFF
QUDA_GPU_ARCH=sm_70
\end{verbatim}

\subsection{colored Galerkin setup 的批量 gather/scatter}

PyQCU 的 strict colored Galerkin 原型原先在每个 column batch 内用
Python 循环逐 source、逐 support 点组装 canonical blocks。该路径在
c128 小格点上也占总 setup 的主要时间，并且放大临时对象数量。本轮保留
每个 color 一次算子调用的账本，但将 blocked-basis gather、fine support
gather 和 block scatter 改为预计算索引的批量张量操作；目标集合和 source
分组分别缓存，仍在实际周期边界上执行冲突检查。

真实 c128 冒烟点 \(8^3\cdot16\) 三层、CPU canonical-block staging、
单精度 coarse blocks 的无缓存 setup 为：\(K=1\) 时 \(44.86\) s，
\(K=4\) 时 \(16.67\) s，约 \(2.69\times\)；两次求解均收敛到
\(1.84\times10^{-8}\) 的最大相对真残差，并保持 35 次外层 FGMRES
迭代。该数字只证明 setup 路径优化，不进入 MG solve 加速比。

进一步的源点 parity 分桶把确定性贪心顺序从 row-major 改为
16 个坐标奇偶类优先。实际周期目标集合仍在桶内做冲突检查，因此正确性
不变，但大格 \(8\cdot16\cdot16\cdot24\) 的源色组由 22 降至 16，
每层 colored operator calls 由 132 降至 96。大格 c128 全 double
C=4/K=16 冷 setup 从约 \(974\) s 降至 \(585.28\) s（transition
分别为 \(351.33\) s 和 \(202.72\) s），已低于同点 QUDA setup
\(630.74\) s；小格 c128 的 setup 也从 \(16.67\) s 降至
\(14.66\) s。

另外，固定输入合成微基准（\(8\cdot16\cdot16\cdot32\)、c128 blocked
basis、c64 blocks、\(K=256\)、44 次算子调用）中，旧 builder 为
\(18.39\) s，新 builder 为 \(5.48\) s，约 \(3.36\times\)；峰值
allocated 均为约 \(13.04\) GiB。该微基准用于隔离 builder 热路径，
不替代真实 gauge 的跨库正式测量。

\section{正式协议}

\subsection{输入与计时}

两侧共享以下输入与边界：

\begin{itemize}
  \item fine Wilson/Clover、mass $0.05$、canonical 12 个 full near-null vectors；
  \item odd--odd MATPC，\code{coarse\_spin=2}，block $[2,2,2,2]$；
  \item c64 tolerance $10^{-6}$、真残差门限 $5\times10^{-6}$；
  \item c128 tolerance $10^{-8}$、真残差门限 $5\times10^{-8}$；
  \item 零初值、2 次 warmup、5 次 measured solve、median/MAD；
  \item steady solve 由调用方自有输入输出计时，setup、输入 I/O、真残差复算分开；
  \item 正式运行设置 \code{QUDA_RESOURCE_PATH} 复用调优参数。
\end{itemize}

\subsection{aligned 与 library 两种 QUDA 策略}

\code{aligned} 将两侧统一到“一个外层迭代执行一次递归 V-cycle”：
MR smoother、\code{smoother\_tol=0}、非最粗层
\code{coarse\_solver\_maxiter=1}，最粗层仍使用 200。

\code{library} 保留 QUDA/PyQUDA 默认：
CA-GCR smoother、\code{smoother\_tol=0.25}、\code{nu\_pre=0}、
\code{nu\_post=8}、\code{coarse\_solver\_maxiter=16}。
该策略只用于判断结论是否依赖某套预算，不与 aligned 数字混合计算。

\section{主结果}

\begin{table}[htbp]
\centering
\small
\caption{aligned 主口径的五次 median；所有正式运行均 trace-off。}
\label{tab:aligned}
\begin{tabular}{@{}L{0.22\textwidth}rrrr@{}}
\toprule
格点 / 层 / 精度 & PyQCU (s) & QUDA (s) & 加速比 & 外迭代 Py/QUDA \\
\midrule
$8^3\cdot16$ / 2 / c64 & 0.054184 & 0.399675 & 7.3763 & 9 / 32 \\
$8^3\cdot16$ / 3 / c64 & 0.043200 & 0.516422 & 11.9542 & 10 / 33 \\
$16^3\cdot16$ / 3 / c64 & 0.149240 & 0.745390 & 4.9946 & 31 / 59 \\
$16\cdot32\cdot32\cdot48$ / 2 / c64 & 1.944663 & 2.183002 & 1.1226 & 11 / 37 \\
$16\cdot32\cdot32\cdot48$ / 3 / c64 & 0.662549 & 1.318356 & 1.9898 & 14 / 39 \\
$8^3\cdot16$ / 3 / c128 & 0.077870 & 0.963774 & 12.3768 & 14 / 46 \\
$16^3\cdot16$ / 3 / c128 & 0.380444 & 1.230048 & 3.2332 & 44 / 87 \\
$16\cdot32\cdot32\cdot48$ / 3 / c128 & 2.065148 & 17.245388 & 8.3507 & 20 / 60 \\
\bottomrule
\end{tabular}
\end{table}

大格 c128 正式点采用 cache-hit：物理缓存由 extended $C=4$、$K=16$、
CPU canonical-block staging 构建，随后按 cache identity 的资产几何与
逐 tensor SHA 校验复用；formal 请求的 build batching 仅作为 provenance，
不参与资产等价判定。该点真残差为 PyQCU $1.78798\times10^{-8}$、QUDA
$1.79244\times10^{-8}$，两侧均通过独立门限。冷 formal $C=1$ setup
尚未完成，因此该 cache-hit 比值不代表冷启动总时间。

表中三层大格 c64、c128 小格和 c128 中格为 builder 优化后的复测值；
与优化前独立运行的对应值相比，MG-vs-MG 加速比仍在约 2--12 倍区间，
变化约在 2--9\% 以内。复测真残差分别为
$4.90135\times10^{-7}$、$1.74532\times10^{-8}$ 和
$1.83242\times10^{-8}$，全部通过正式门限。

小格点加速比明显大于中型/大型格点，原因是固定 setup 之外的外层迭代数
差异在小体积上占主导；大格点则同时受每层算符规模、缓存与固定 kernel
开销影响。不能仅凭单一小格点外推正式部署性能。

同一 $16\cdot32\cdot32\cdot48$ c64 格点上，二层为
$1.1226\times$，三层为 $1.9898\times$。这不是单点偶然值，而是层数
消融：增加第三层后 PyQCU 时间从 $1.944663$ s 降至 $0.662549$ s，
QUDA 从 $2.183002$ s 降至 $1.318356$ s，两侧均从更深粗化中获益，
但 PyQCU 的降幅更大。

\begin{table}[htbp]
\centering
\small
\caption{QUDA 原生 library 策略敏感性；不替代表 \ref{tab:aligned} 主口径。}
\label{tab:library}
\begin{tabular}{@{}L{0.28\textwidth}rrr@{}}
\toprule
格点 / 层 / 精度 & PyQCU (s) & QUDA library (s) & 加速比 \\
\midrule
$16^3\cdot16$ / 3 / c64 & 0.170131 & 0.465514 & 2.7362 \\
$16\cdot32\cdot32\cdot48$ / 3 / c64 & 0.654533 & 2.462287 & 3.7619 \\
$16^3\cdot16$ / 3 / c128 & 0.357061 & 0.961293 & 2.6922 \\
\bottomrule
\end{tabular}
\end{table}

\section{BiCGStab 参考与失真检查}

同侧 BiCGStab 不进入 MG-vs-MG 加速比，但用于识别计时或配置失真。

\begin{table}[htbp]
\centering
\small
\caption{三层复测：MG 与同侧 plain BiCGStab median。}
\label{tab:bicg}
\begin{tabular}{@{}L{0.23\textwidth}rr@{}}
\toprule
实现 & MG (s) & BiCGStab (s) \\
\midrule
PyQCU & 0.662549 & 1.685399 \\
QUDA & 1.318356 & 0.325375 \\
\bottomrule
\end{tabular}
\end{table}

表 \ref{tab:bicg} 为大格 c64；同一协议的 c128 小格/中格复测为
PyQCU MG/BiCGStab $0.077870/0.304677$ s，QUDA
$0.963774/0.103475$ s（小格）；PyQCU
$0.380444/0.453031$ s，QUDA $1.230048/0.118572$ s（中格）。
大格 c128 为 PyQCU MG/BiCGStab $2.065148/3.229977$ s，QUDA
$17.245388/0.978754$ s。
因此只有 MG-vs-MG 比值可以归因于 MultiGrid 路径：例如 c128 中格的
$3.23\times$ MG 比值并不表示 PyQCU 的总体求解器快于 QUDA，因为 QUDA
BiCGStab 仍比 PyQCU MG 快约 $3.2\times$，也比 PyQCU BiCGStab 快约
$3.8\times$。大格 c64 上 QUDA BiCGStab 比 PyQCU MG 快约 $2.0\times$。
正式结论必须同时报告这组参考时间，不能把 MG 路径优势扩大解释为一般
求解器优势。

此外，c128 小格采用单精度 coarse blocks 并把 Galerkin batch 调到
$K=4$ 时，PyQCU setup 从 $44.86$ s 降至 $16.67$ s，但外层迭代由
double coarse blocks 的 14 次增加到 35 次，steady solve 为
$0.177628$ s；该结果是显存/容量实验，不是与 double-coarse QUDA 的
主口径等精度比较。

\section{trace 开销}

小格 c64 三层同一参数下，trace-off 取正式文档的 median，trace-on 取
诊断文档的 steady median。trace-on 数据只用于量化观测代价，不参与加速比。

\begin{table}[htbp]
\centering
\small
\caption{trace-off 与 trace-on 的 steady median。}
\label{tab:trace}
\begin{tabular}{@{}L{0.25\textwidth}rrr@{}}
\toprule
模式 & PyQCU (s) & QUDA (s) & 说明 \\
\midrule
trace-off & 0.043971 & 0.505706 & 正式主口径 \\
trace-on & 0.125649 & 0.560432 & 诊断，excluded \\
开销因子 & $2.86\times$ & $1.11\times$ & 比较同侧 \\
\bottomrule
\end{tabular}
\end{table}

PyQCU 的 trace 需要在 C++ 内记录逐层残差与阶段，并执行额外同步，因此
对微小格点的相对开销远高于 QUDA 的原生日志；这也验证了正式速度比必须
在 trace-off 下产生。

\subsection{大格 c128 的 PyQCU 分层诊断}

在 formal cache-hit 的大格 c128 上再次开启 PyQCU trace，8 次 solve
（2 warmup、5 measured、1 untimed memory probe）均保持 20 次外迭代，
trace-on 墙钟约 $3.31$--$3.43$ s，而 trace-off steady 为
$2.065148$ s；观测开销约 $1.6\times$。阶段计时中，单次 solve 的
outer iteration 合计约 $3.374$ s，其中 coarse V-cycle 约 $2.146$ s
（占 outer 约 $63.6\%$）；fine pre/post smoother 合计约 $0.248$ s，
fine restriction/prolongation 合计约 $0.247$ s，fine correction
residual 约 $0.112$ s。最粗层 BiCGStab 在单次 solve 中约 $0.297$ s。

该 trace 说明 PyQCU 的 $8.3507\times$ MG 优势不能仅归因于某个单层
kernel：每外层迭代的 coarsest recursion 和 coarse V-cycle 是主要成本，
但外层迭代数明显少于 QUDA 的 60 次。由于本轮没有为大格 QUDA 重跑
trace，且 trace 本身有约 $1.6\times$ 开销，这些数字只用于解释算法
结构，不能替换 trace-off 正式结果。

\section{正确性与残差}

所有正式点均通过独立全 Wilson/Clover 真残差复算。大格 c64 三层的
最大相对真残差为 PyQCU $4.90135\times10^{-7}$、QUDA
$7.32826\times10^{-7}$；中型 c128 三层分别为
$1.83242\times10^{-8}$ 与 $1.83035\times10^{-8}$。
两库输出均来自同一 RHS，且解在进入残差复算前按约定完成归一化转换。

\section{显存与多卡边界}

\subsection{c128 大格点}

c128 三层 $16\cdot32\cdot32\cdot48$ 的早期 $C=1$ 冷 setup 在
32 GiB V100 上 OOM 或超过 30 分钟；现在的 extended builder 使用
$C=4$、$K=16$、CPU canonical-block staging、null-basis offload，
完成后将同一物理缓存提供给 formal $C=1$ solve，并在 runtime seal 后
调用 \code{empty\_cache()} 归还 setup 临时块。PyQCU 两侧缓存构建
setup 统计为 $351.33$ s 与 $202.72$ s（总 setup $585.28$ s），
QUDA double setup 为 $630.74$ s。formal solve 的 MG steady 为
PyQCU $2.065148$ s、
QUDA $17.245388$ s，正式比值 $8.3507\times$；相对真残差分别为
$1.78798\times10^{-8}$ 与 $1.79244\times10^{-8}$，均通过独立门限。
冷 formal $C=1$ setup 仍未完成，因此不能把 cache-hit solve 比值当作
冷启动总时间比值。

同侧 plain BiCGStab 为 PyQCU $3.22998$ s、QUDA $0.97875$ s；
QUDA BiCGStab 比 PyQCU MG 快约 $2.11\times$，也比 PyQCU BiCGStab
快约 $3.30\times$。因此该 $8.3507\times$ 只能解释为两条 MultiGrid
路径的差异，不能外推为 PyQCU 一般求解器优势。

\subsection{QUDA mixed-coarse 容量补充}

为区分“QUDA 算法不可用”与“全 double 容量不足”，另运行了
\code{--quda-coarse-precision single} 的 QUDA-only 点：
外层仍为 c128，QUDA coarse MG 使用 single precision，gate 仍由独立
c128 真残差复算。$16\cdot32\cdot32\cdot48$ 三层成功，setup 约
$617$ s，60 次外层迭代的 steady solve 为 $2.265$ s，最终相对真残差
$1.79\times10^{-8}$；运行期间设备峰值约 $9.8$ GiB。

该点证明了 QUDA 的 c128 mixed-coarse 路径可用并显著降低显存占用，但
PyQCU strict 当前显式要求各层相同 dtype，尚未实现同等 mixed-coarse
runtime，因此不能把 QUDA-only 数字写入跨库 speedup。PyQCU 的
block-precision 低内存选项虽把 coarse blocks 降为单精度，仍在 batch
diagonal inverse 阶段 OOM；进一步改用分块 site inverse 并将 projection
batch 降到 1 后，可越过该点；随后又把分块展平到 site 轴并使用 compact
complex64 inverse cache，最终仍在 fine dslash batch matvec 处 OOM。这
作为后续 strict mixed-precision 实现的入口，而不是已完成结论。
失败文档保留在 \pathref{data/mg_matrix_20260916_round2/formal-c128-l3/benchmark.json}。
因此 c128 最大可信跨库正式点是 $16^3\cdot16$ 三层；小格 c128 三层
用于确认 double coarse MG 的真实 setup/solve，而不是用单侧结果伪造比值。

\subsection{多卡}

当前运行环境只有一张 V100 属于 PyTorch 构建支持的 \code{sm_70+} 架构。
两张 P100 为 \code{sm_60}，当前 PyTorch 支持列表不含 \code{sm_60}，
QUDA 本轮也按 \code{sm_70} 构建。因此不能把 P100 并发纳入本轮正式
MultiGrid 结论；这属于硬件/工具链能力缺口，不是算法验证失败。阶段 1
MPI 能力门禁仍在可用范围内通过：单 rank preflight 为
$18$ passed/$9$ skipped，\code{mpirun -np 2} 每个 rank 为
$14$ passed/$13$ skipped；这些测试只确认 c64/c128 全局 dot/norm
归约与 fail-closed preflight，不能解释为已实现分布式 halo 或
distributed fused FGMRES。

\section{复现命令与证据}

\begin{verbatim}
source ./env.sh
source examples/qcu/dev87/quda_env.sh
export QUDA_RESOURCE_PATH=/root/PyQCU/data/quda-resource-cache
export QUDA_ENABLE_TUNING=1
unset PYQCU_STRICT_TRACE_FILE QUDA_MG_TRACE_FILE

# 主口径：正式大格 c64 三层
python -B examples/qcu/dev87/bench_strict_vs_quda.py \
  --profile formal --side both --quda-strategy aligned \
  --lattice 16 32 32 48 --levels 3 --precision c64 \
  --repeats 5 --cache-expect hit \
  --reference-solver bicgstab --reference-warmups 1 \
  --gauge-path data/gauge_16x32x32x48_m0.05_seed42_c64.h5 \
  --nullvec-path data/L16x32x32x48_nvec12_full_c64.h5 \
  --quda-nullvec-prefix data/L16x32x32x48_nvec12_quda \
  --quda-nullvec-manifest data/L16x32x32x48_nvec12_quda.conversion.json \
  --output data/mg_matrix_20260916_round2/formal-c64-l3/benchmark.json

# trace 诊断：正式加速比之外单独报告
python -B examples/qcu/dev87/trace_strict_vs_quda.py \
  --profile smoke --repeats 2 --lattice 8 8 8 16 --levels 3 \
  --precision c64 --cache-expect any \
  --reference data/mg_matrix_20260916_round2/trace-reference-small-c64-l3.json \
  --output data/mg_matrix_20260916_round2/trace-small-c64-l3.json
\end{verbatim}

\begin{longtable}{@{}L{0.34\textwidth}L{0.60\textwidth}@{}}
\caption{主要证据文件}\label{tab:evidence}\\
\toprule 文件 & 内容\\\midrule
\pathref{data/mg_matrix_20260916_round2/formal-c64-l3/benchmark-main-recheck.json} &
当前 c64 大格三层 aligned 五次正式复测记录。\\
\pathref{data/mg_matrix_20260916_round2/formal-c64-l3-medium/benchmark.json} &
c64 中格三层 formal。\\
\pathref{data/mg_matrix_20260916_round2/formal-c64-l2-small/benchmark.json} &
c64 小格二层 formal。\\
\pathref{data/mg_matrix_20260916_round2/formal-c64-l3-small/benchmark.json} &
c64 小格三层 formal。\\
\pathref{data/mg_matrix_20260916_round2/formal-c128-l3-small/benchmark-main-recheck.json} &
c128 小格三层 formal 复测，double coarse 主口径。\\
\pathref{data/mg_matrix_20260916_round2/formal-c128-l3-small/benchmark-k4.json} &
c128 小格三层 single coarse blocks/CPU staging 容量实验；不进入主 speedup。\\
\pathref{data/mg_matrix_20260916_round2/formal-c128-l3-medium/benchmark-main-recheck.json} &
c128 中格三层 formal 复测，double coarse 主口径。\\
\pathref{data/mg_matrix_20260916_round2/formal-c128-l3/combined-double-c4-k16-smoke.json} &
c128 大格三层全 double 双侧 smoke；$C=4$、$K=16$、CPU staging，
不进入 formal speedup。\\
\pathref{data/mg_matrix_20260916_round2/formal-c128-l3/benchmark-formal.json} &
c128 大格三层 formal cache-hit；缓存由 extended $C=4$ 构建并经资产
identity 校验，主 MG 比值为 $8.3507\times$。\\
\pathref{data/mg_matrix_20260916_round2/formal-c128-l3/pyqcu-double-c4-k16-parity.json} &
c128 大格 parity 分组后的冷 C=4 全 double setup；总 setup $585.28$ s，
solve 真残差通过。\\
\pathref{data/mg_matrix_20260916_round2/formal-c128-l3/pyqcu-formal-trace.tsv} &
大格 c128 PyQCU trace-on 的逐层 stage/residual 诊断；不进入正式计时。\\
\pathref{data/mg_matrix_20260916_round2/library-formal-c64-l3/benchmark.json} &
QUDA library 策略敏感性。\\
\pathref{data/mg_matrix_20260916_round2/trace-small-c64-l3.json} &
trace-on/off 小格诊断与残差阶段。\\
\bottomrule
\end{longtable}

\section{结论}

本轮得到三点可复核结论：

\begin{enumerate}
  \item 旧 c64 $4.3817\times$ 来自未对齐内层预算，已撤回；当前大格 c64
        三层 aligned 主结果为 $1.9898\times$，五次样本和真残差均通过。
  \item QUDA double MG 与 PyQCU double MG 均已推进到 c128 小/中格点
        以及大格三层；大格 formal cache-hit 的 MG 比值为
        $8.3507\times$，但冷 $C=1$ setup 仍未完成，且 QUDA plain
        BiCGStab 仍快于 PyQCU MG，故不能外推为一般求解器优势。
  \item trace 功能对 PyQCU 小格点有显著观测开销，正式加速比必须在
        trace-off 下采集；BiCGStab 参考和 aligned/library 双策略已作为
        后续防失真硬闸门写入技能库。
\end{enumerate}

\end{document}
